\section*{Chapter \ref{ch:m1:margin} --- Margin}

\begin{solution}{exo:m1:margin:1}
$-44.25$ points $\times\ \$50 \times 25 = -\$55\,312.50$: the account pays.
\end{solution}

\begin{solution}{exo:m1:margin:2}
Equity falls to $120\,000 - 16\,000 = \$104\,000$, below the maintenance
requirement of $5 \times 21\,000 = \$105\,000$: a call, to restore the initial
level of $5 \times 23\,100 = \$115\,500$. The client must deposit \$11\,500.
\end{solution}

\begin{solution}{exo:m1:margin:3}
One third of the range is 115 points, \$5\,750. Scenarios 1--2: 0, 0. Then
(loss positive) 3--4: $-5\,750$; 5--6: $+5\,750$; 7--8: $-11\,500$; 9--10:
$+11\,500$; 11--12: $-17\,250$; 13--14: $+17\,250$; 15: $-15\,525$; 16:
$+15\,525$ ($3 \times 17\,250 \times 30\%$). The margin is the largest loss:
\$17\,250, scenarios 13 and 14.
\end{solution}

\begin{solution}{exo:m1:margin:4}
The portfolio is short options, hence short volatility: every odd scenario
adds four volatility points and raises the value of what it has sold.
Scenario 16 is worse than 15 because the portfolio is also long 2 futures:
a fall hurts the futures and the short puts together, while in a rise the
futures' gain offsets part of the short calls' loss. Scenario 13 (full range
down, volatility up) would bind if the options were struck closer to the
money, so that most of their loss occurs within one scan range, or if the
extreme scenarios counted a smaller fraction.
\end{solution}

\begin{solution}{exo:m1:margin:5}
$213\,250 - 107\,770 = \$105\,480$, 49\%. At 5\% a year: \$5\,274 a year for this
small portfolio, and, more importantly, half the position size for the same
capital.
\end{solution}

\begin{solution}{exo:m1:margin:6}
$2.1\% \times \sqrt2 = 2.97\%$; $2.1\% \times \sqrt5 = 4.70\%$. Higher: volatility
clusters and returns in a crisis are positively autocorrelated, and the
liquidation itself moves prices. Lower: mean reversion in some products, and
a portfolio that is partly hedged on day one does not carry its full risk for
five days.
\end{solution}

\begin{solution}{exo:m1:margin:7}
Ten-day increases: 109\%, 203\% and 94\%. Calm averages: 3.0\% (filtered) and
4.5\% (floored at a weight of 25\%). A weight of 47\% brings the ten-day
increase to 50\%, at a calm-period margin of 5.8\% of the position, nearly
double the unfloored figure: stability is bought with collateral that sits
idle most of the time.
\end{solution}

\begin{solution}{exo:m1:margin:8}
In a 10\% \emph{rise} the futures lose \$20 million, payable in cash the next
morning, while the \$20 million gained on the shares is unrealised: the \omterm{def:m1:financing:pb}{prime
broker} will lend against it, but at a haircut, on its own timetable, and not
at all if it is cutting leverage at the same moment. In the fall the cash
arrives from the clearing house, but the \omterm{def:m1:financing:pb}{prime broker} may call margin on the
shares. In both directions \omterm{def:m1:clearing-and-settlement:margin}{initial margin} on the futures rises with
volatility. Market neutral is a statement about P\&L; margin is a statement
about cash, leg by leg, venue by venue.
\end{solution}

\begin{solution}{pb:m1:margin:1}
\textbf{1.} $400 \times 3\,300 \times 50 = \$66$ million; margin $4.5\% =
\$2.97$ million.
\textbf{2.} $400 \times 50 \times 800 = \$16$ million.
\textbf{3.} $400 \times 2\,500 \times 50 \times 9\% = \$4.5$ million: $+\$1.53$
million, although the notional has fallen by a quarter.
\textbf{4.} \$17.53 million, 26.6\% of the initial notional.
\textbf{5.} It has \$12 million. It must raise \$5.5 million by selling
bonds, in the month when they are hardest to sell, or cut the position at
the lows.
\textbf{6.} The future falls to 2\,488.75: $400 \times 50 \times 261.25 =
\$5.225$ million.
\textbf{7.} $400 \times 50 \times 165 = \$3.3$ million, due within the hour:
cash that is in a bond fund settling in two days does not count. A missed
intraday call is a default like any other.
\textbf{8.} From $6\% \times 55$ million $= \$3.3$ million to $7.5\% \times
49.775$ million $= \$3.733$ million: $+\$433\,125$.
\textbf{9.} $5\,225\,000 + 433\,125 = \$5\,658\,125$, of which \$3.3 million at
noon.
\textbf{10.} The filtered model the most (margin tripled in ten days); the
floored model the least in relative terms, the plain historical model the
least in absolute terms because it barely moved.
\textbf{11.} The filtered and the floored models, which by then are near
9\% of the position; the plain historical window is still near 5\% while
daily volatility is 3.5\%, so a two-day move at 99\% is about 11\%: it is
under-margined by half.
\textbf{12.} It asks for half as much again in calm years, when members see
no risk, competitors advertise lower margins, and the cost is visible while
the benefit is not.
\textbf{13.} The second, if the fund sizes its positions to it: its calls in
the crisis are proportionally the same, but the fund has been holding less
leverage and more collateral all along. The first is cheaper every day but
one.
\textbf{14.} Every long and every short posts \omterm{def:m1:clearing-and-settlement:margin}{initial margin}; doubling the
rate doubles the collateral locked at the clearing house on both sides of
every contract, while \omterm{def:m1:clearing-and-settlement:margin}{variation margin} merely moves cash from losers to
winners.
\textbf{15.} Not with the cash. The shares' gain is at a \omterm{def:m1:financing:pb}{prime broker}, which
decides how much can be withdrawn; the futures' loss is due in cash at the
clearing house in the morning. The hedge removes the loss, not the call.
\textbf{16.} The cash needed to hold the position through a stress: a month
of adverse \omterm{def:m1:clearing-and-settlement:margin}{variation margin} at a stressed volatility, plus the increase of
\omterm{def:m1:clearing-and-settlement:margin}{initial margin} under the clearing house's model in the same scenario,
against cash and assets that are actually eligible and available in hours.
\textbf{17.} Clearing houses accept a list of collateral, mostly cash and
government bonds, with haircuts; corporate bonds are often ineligible or
heavily haircut, and \omterm{def:m1:clearing-and-settlement:margin}{variation margin} is cash only.
\textbf{18.} Bond prices fall, funds holding them face redemptions and their
own calls, volatility rises, and the models of
\cref{fig:m1:margin:procyclical} raise margins again: the loop that the
word \omterm{def:m1:margin:procyclical}{procyclicality} names.
\textbf{19.} \textbf{\$1.53 million.}
\textbf{20.} Solvency is about the value of assets against liabilities;
a \omterm{def:m1:financing:margincall}{margin call} is about cash in the right currency at the right bank by ten
o'clock.
\end{solution}

\begin{solution}{iq:m1:margin:1}
\omterm{def:m1:clearing-and-settlement:margin}{Variation margin} settles the day's gain or loss in cash, so that no debt
accumulates between the member and the clearing house; it passes from losers
to winners. \omterm{def:m1:clearing-and-settlement:margin}{Initial margin} is collateral held against the loss that could
occur between a default and the close-out of the defaulter's positions; it
is returned when the position is closed. The first is P\&L, the second a
deposit whose size the clearing house can change.
\iqlookfor{the purpose of each, and that \omterm{def:m1:clearing-and-settlement:margin}{initial margin} can change with no
change in position.}
\end{solution}

\begin{solution}{iq:m1:margin:2}
For each product the clearing house sets a price and a volatility scan
range. Every contract has a risk array: its loss under sixteen scenarios,
price moves of thirds of the range with volatility up and down, plus two
extreme moves counted in part. Multiply by positions, sum by scenario, take
the worst: the \omterm{def:m1:margin:scan}{scan risk}. Add a charge for spreads between expiries, apply a
minimum per short option, subtract credits between correlated products.
\iqlookfor{the extreme scenarios and the reason for the short option
minimum.}
\end{solution}

\begin{solution}{iq:m1:margin:3}
Margins that rise in stress add to the demand for cash when cash is
scarcest. Tools: a floor from a stressed period, a weight on stressed
value-at-risk, a long lookback that always contains a crisis, a buffer built
in calm times and released in stress, limits on the speed of increases. All
of them raise the average margin, which members pay for every day, and the
last one leaves the clearing house under-covered while it catches up.
\iqlookfor{at least two tools and the cost of each.}
\end{solution}

\begin{solution}{iq:m1:margin:4}
The extreme scenarios and the short option minimum. Within the normal scan
range a far out-of-the-money option loses almost nothing, so a scan alone
would let a seller accumulate unlimited size for almost no margin, a
position that loses catastrophically in exactly the move the scan does not
reach. The extreme scenarios revalue at two or three ranges and count a
fraction of the loss; the minimum puts a floor under every short option.
\iqlookfor{why the standard scenarios are blind to the position.}
\end{solution}

\begin{solution}{iq:m1:margin:5}
Track cash, not only P\&L: daily \omterm{def:m1:clearing-and-settlement:margin}{variation margin}, \omterm{def:m1:clearing-and-settlement:margin}{initial margin} under the
clearing house's method with its parameters as they were on each date
(margin rates rise in crises, precisely when a leveraged strategy is
losing), the broker's multiplier, and interest on posted collateral. Impose
the constraint that cash never goes negative; size positions to the stressed
margin, not the current one. Report the return on the capital actually
needed, including the buffer.
\iqlookfor{historical margin rates and the cash constraint.}
\end{solution}

\begin{solution}{iq:m1:margin:6}
P\&L nets across the two houses; cash does not. On any large move one house
calls \omterm{def:m1:clearing-and-settlement:margin}{variation margin} on the losing leg in the morning while the gain at
the other is paid out on that house's cycle, possibly later and in another
currency. \omterm{def:m1:clearing-and-settlement:margin}{Initial margin} rises at both. Size the buffer from a joint stress:
the largest multi-day move in the common risk factor, with each house's
model re-run at stressed volatility, plus the timing gap between paying one
and receiving from the other, plus currency. Pre-arrange credit lines and
collateral transformation; test the intraday call operationally.
\iqlookfor{the timing gap and the simultaneous rise in \omterm{def:m1:clearing-and-settlement:margin}{initial margin}.}
\end{solution}
